\section{Literature Review}
\label{sec:literature-review}

Researchers have studied ways to solve the ride-sharing problem, as well as matching and scheduling, extensively over the last 20 years. The past literature approaches the ride-sharing problem from both operations research and systems perspectives. This section reviews the history of the ride-sharing problem, how it has been formulated and solved, and the history of matching and scheduling in the literature. 

The operations research view defines the problem as a routing problem and places it in the category of pickup-and-delivery problems, in which vehicles transport loads from origins to destinations~\citep{savelsbergh1995general}. The dial-a-ride problem is the special case where the loads are people~\citep{savelsbergh1995general}, ride-sharing is similar to the dynamic variant of the pickup-and-delivery problem~\citep{furuhata2013ridesharing}, and the customer experience is included in the objective~\citep{cordeau2007dialaride}. \citet{psaraftis1980dynamic} gives an exact dynamic programming algorithm for a single-vehicle many-to-many problem that minimizes a weighted combination of total service time and customer dissatisfaction. \citet{psaraftis1980dynamic} also shows that computational effort grows exponentially with problem size. Exact algorithms remain limited to small instances, while heuristics are the only methods that reach large ones~\citep{cordeau2007dialaride}. Also, in the dynamic case, where requests are revealed over time, a solution can no longer be a fixed plan but a strategy that reacts to new information~\citep{berbeglia2010dynamic}. The matching part of the problem has an even longer history as an assignment problem~\citep{pentico2007assignment}. Quite recently, \citet{masoud2017realtime} match each customer one at a time against all drivers, allowing riders to transfer between vehicles, and solve this many-to-one problem optimally with a dynamic program over a time-expanded network.

The literature that views the problem from a systems perspective investigates and builds matching programs to solve it. \citet{hall1997dynamic} studied dynamic ride-sharing, both analytically and through an implementation in Los Angeles, and found that a busy road produces a large enough pool of potential matches in theory. In practice, however, a customer had, at best, a one-in-five chance of being offered a ride per person contacted, even in an emergency. Hence, a match that is feasible in theory is not necessarily a match people accept~\citep{hall1997dynamic}. \citet{dailey1999seattle} moved matching to a series of web page forms, ran the system alongside a traditional phone-based program for a year, and reached about the same number of new users, suggesting that the matching mechanism does not limit it. Their model makes the number of matches proportional to the probability that two trips are compatible, so any restriction on a match lowers the number of matches formed~\citep{dailey1999seattle}. \citet{agatz2012optimization} survey the optimization side of the dynamic problem and conclude that centralized matching may be computationally infeasible for large cities, so a decomposition approach is necessary.

\subsection{Ride-sharing}

The term "ride-sharing" is not defined consistently in the literature. The broadest definition, by \citet{furuhata2013ridesharing}, defines it as a joint trip in which more than one participant shares a vehicle. \citet{furuhata2013ridesharing} also require participants to coordinate where and when each person is picked up and dropped off. \citet{chan2012ridesharing} add that the driver shares the same origin or destination as the passengers and does not financially benefit from the arrangement. \citet{agatz2012optimization} give a similar definition but also require that drivers are "independent private entities" and that trips are "single, non-recurring". \citet{mitropoulos2021systematic} find that, in the literature, the term is used for "both profit and non-profit ride-sharing services". \citet{pigalle2023ridesharing} separate ride-sharing from ride-hailing, which refers to the for-profit connection between riders and drivers. The organized variants share prearrangement by a service provider, which separates ride-sharing from hailing a taxi on the street~\citep{furuhata2013ridesharing}.

The literature is also divided by who runs the ride-sharing service. \citet{chan2012ridesharing} classify ride-sharing by the relationships among participants, separating acquaintance-based arrangements from organization-based ones where participants have a membership or join through a website. \citet{furuhata2013ridesharing} make a similar distinction on the provider side, distinguishing matching agencies that only introduce independent drivers and customers and service operators that have their own vehicles and drivers. These two cases have different problems. For an agency, both sides must screen the other side themselves, making trust a precondition for matching~\citep{chaube2010trust}. For an operator, they make most decisions, while the customer only needs to decide whether to take part~\citep{furuhata2013ridesharing}. This gives the fixed-fleet problem familiar from operations research, where vehicles have capacities and fixed start and end locations \comments{I coulnd't understand this sentence.}~\citep{savelsbergh1995general}, and requests can be turned down when the drivers cannot serve them~\citep{cordeau2007dialaride}. The scenario studied in this thesis falls into the operator category, where drivers work fixed shifts.

\subsection{Matching}

Matching customers to drivers is, in its most general form, an assignment problem, another category of problems which has a long history and many variants and algorithms~\citep{pentico2007assignment}. A classic assignment problem makes one-to-one assignments to minimize total cost, a problem that can be solved efficiently~\citep{pentico2007assignment, edmonds1972theoretical}. Ride-sharing differs in a few ways. The assignment is usually not a one-to-one matching, as a driver most likely serves several customers during a shift~\citep{savelsbergh1995general} and, when a matching agency does the matching, a customer is offered more than one driver to choose from~\citep{dailey1999seattle}. Also, the cost of an individual match is unknown before assignment because it depends on the driver's other matches. In the dial-a-ride problem, the feasibility of adding a request depends on the vehicle's current route~\citep{cordeau2007dialaride}. The objective can then be the number of matches rather than total cost, which \citet{agatz2012optimization} surveyed as an objective. Such a capacitated bipartite matching can be expressed as a maximum flow problem, in which the number of drivers a customer may be offered and the number of customers a driver may serve become edge capacities~\citep{ford1957simple}, and it is solvable in polynomial time with an integral solution~\citep{edmonds1972theoretical}.

\citet{masoud2017realtime} state that matching is integral to ride-sharing, and they incorporate time feasibility by matching a customer only when a feasible schedule exists for the ride. This avoids the poor scaling of exact dial-a-ride algorithms~\citep{cordeau2007dialaride}. The decision-maker in the matching also matters: when a service operator does the matching, it makes most decisions, whereas a matching agency introduces counterparts and leaves the decision to them \comments{Who's them here? It is a bit ambiguous}~\citep{furuhata2013ridesharing}. Under an operator, customer preferences can therefore enter the matching only as constraints. The preferences follow from the relationship between participants and the trust between the two sides~\citep{chaube2010trust}. Ride-sharing is sparse in time and space, which limits the number of customers that can be served~\citep{masoud2017realtime}, and the more preferences a customer has, the harder it is to find a suitable match~\citep{dailey1999seattle,agatz2012optimization}. This is consistent with the earlier observations that people do not always accept a match based on theory~\citep{hall1997dynamic}.

\subsection{Scheduling}

Choosing when each driver's ride takes place is a scheduling problem in the classical sense. In machine scheduling, each job has a processing time, a weight, a release date when the job becomes available, and a due date, and each machine can handle at most one job at a time~\citep{lenstra1977complexity}. In these terms, a driver is a single machine, a ride is a job whose processing time is the ride duration, and the wanted pickup time is the release date. The optimality criteria are functions of job completion times, such as maximum lateness, total tardiness, or the number of late jobs~\citep{lenstra1977complexity}. A late pickup is then a tardiness, and an unserved request equals a late job. \comments{Release dates alone make a single-machine problem hard. Not sure what this means.}. \citet{lenstra1977complexity} show that a non-zero release date for a single job is enough to make minimizing maximum lateness, total completion time, or the number of late jobs NP-complete \comments{It took a bit but I understood, but let us re-write this a bit}. Therefore, exact solutions rely on integer programming. \citet{dyer1990formulating} compare several formulations of the single-machine problem with release dates and conclude that formulations based on time discretization give the best lower bounds. In their best formulation, a binary variable equals one if a job starts processing at a given time, which makes the feasible set essentially independent of the data but gives the formulation a very large number of variables~\citep{dyer1990formulating}. However, classical criteria do not measure the machine, since they are all functions of job completion times~\citep{lenstra1977complexity} \comments{Might need an extra explanation}. In the dial-a-ride problem, the vehicle's waiting time at a stop is a scheduling variable, and some variants bound it from above, but drivers remain in the objective only through their wages~\citep{cordeau2007dialaride}.

In the dial-a-ride literature, the scheduling step follows the routing step~\citep{cordeau2007dialaride,savelsbergh1995general}. Given a route, the scheduling problem determines the departure time from the depot and the time the ride begins at each stop so that the time constraints are satisfied and the route duration is minimized~\citep{cordeau2007dialaride}. For a fixed route, this is altrivial, as the optimal departure time can be determined in time linear in the length of the route~\citep{cordeau2007dialaride}. The difficulty comes from choosing the route, since for a single vehicle the exact dynamic program of \citet{psaraftis1980dynamic} grows exponentially in the number of requests. Multi-vehicle methods therefore decompose the problem~\citep{savelsbergh1995general,cordeau2007dialaride} \comments{I think that something is missing here}. A common technique defines customer clusters to be served by the same driver before the routing phase, and then solves each cluster using a single-vehicle algorithm~\citep{cordeau2007dialaride,savelsbergh1995general}. Some other methods alternate between a routing step and a scheduling step~\citep{cordeau2007dialaride,savelsbergh1995general}. Once a driver's customers are assigned and each ride is defined as a block of time, we have the single-machine problem described above for each vehicle. The clusters are not independent, however, and early cluster-first \comments{Let us use a different word here} heuristics already switch between clusters after solving each one~\citep{cordeau2007dialaride}. \citet{agatz2012optimization} note that a geographic partition can be problematic because a ride's origin and destination may fall in different subregions. \citet{masoud2017realtime} take the opposite route and keep time feasibility inside the matching, constructing a time-expanded feasible network for one rider at a time and searching it with a dynamic program, at the cost of processing the riders one after another.

\subsection{Contribution}

This thesis investigates the problem by decomposing it into matching, per-driver scheduling, and conflict resolution, following the decomposition that \citet{agatz2012optimization} expect to be necessary. Matching on the preference graph alone is inexpensive, since maximum flow is solvable in polynomial time~\citep{edmonds1972theoretical}, and it reduces each driver's candidate set before scheduling. The conflict resolution step plays the role of the swaps between clusters that the cluster-first \comments{same here} heuristics make after solving each cluster~\citep{cordeau2007dialaride}.

We add customer preferences as hard constraints to the matching instead of relying on users to screen for their driver. The customers' preferences form a bipartite graph in which suitable drivers are matched to each customer. In practice, screening is done by manual selection from stored participant profiles~\citep{furuhata2013ridesharing}; whether a match is accepted depends on trust between the parties~\citep{chaube2010trust}; and restrictions are treated as something that makes matching harder~\citep{agatz2012optimization}. \citet{furuhata2013ridesharing} identify integrating screening into the ride-matching itself, which they call high-dimensional matching, as an open challenge, and note that screening items treated as hard constraints are the simple case. This thesis takes that simple case: each customer's list of acceptable drivers is a hard constraint on the matching.

Past literature also rarely accounts for drivers when forming model objectives, and it often uses decomposition. Drivers usually appear only as wage costs~\citep {cordeau2007dialaride}. Instead, we add driver idle time as a term in the objective and try to minimize it. Arguably, we should also account for driver experience, which we investigate here.